% Generated draft: requires mathematical and source review.
\documentclass[11pt]{article}
\usepackage[margin=25mm]{geometry}
\usepackage{fontspec}
\setmainfont{Latin Modern Roman}
\usepackage{amsmath,amssymb,mathtools,unicode-math}
\setmathfont{Latin Modern Math}
\usepackage{enumitem,needspace,xurl,hyperref}
\setlength{\parindent}{0pt}
\setlength{\parskip}{5pt}
\setlength{\emergencystretch}{3em}
\newcommand{\N}{\mathbb{N}}
\newcommand{\Z}{\mathbb{Z}}
\newcommand{\Q}{\mathbb{Q}}
\newcommand{\R}{\mathbb{R}}
\newcommand{\C}{\mathbb{C}}
\newcommand{\F}{\mathbb{F}}
\newcommand{\A}{\mathbb{A}}
\newcommand{\PP}{\mathbb{P}}
\newcommand{\E}{\mathbb{E}}
\newcommand{\eps}{\varepsilon}
\newcommand{\dd}{\,\mathrm{d}}
\begin{document}
\section*{UnsolvedMath ID: 800923}
\textbf{Source problem code:} catalogue entry
\subsection{Definitions and mathematical statement}
Let \(F\) be a smooth, strongly convex Finsler metric on \(S^2\); reversibility is not assumed. Thus \(F:TS^2\to[0,\infty)\) is smooth off the zero section, positive on nonzero vectors, positively homogeneous of degree one, and the fibrewise Hessian of \(\tfrac12F^2\) is positive definite off the zero section. A closed geodesic is a nonconstant periodic geodesic, parametrized with unit speed. It is \emph{prime} if it is not an iterate of a closed geodesic of smaller period. Write \(\mathcal P(F)\) for the set of prime closed geodesics modulo translation of the unit-speed parameter.

For \(c\in\mathcal P(F)\), let \(P_c\) be the linearized Poincar\'e return map on a two-dimensional transverse section in the unit tangent bundle. Call \(c\) \emph{hyperbolic} if
\[
\operatorname{spec}(P_c)\cap\{z\in\C:|z|=1\}=\varnothing.
\]
Under the non-strict spectral convention, call \(c\) \emph{elliptic} if
\[
\operatorname{spec}(P_c)\subseteq\{z\in\C:|z|=1\}.
\]
This convention permits eigenvalues \(1\) and \(-1\), including degenerate cases. The supplied source does not specify whether these cases are intended to count as elliptic; this qualification remains unresolved.

The supplied record contains the following two conjectural assertions, with the displayed conjecture attributed to Long:
\begin{enumerate}
\item For every Finsler metric \(F\) on \(S^2\),
\[
\bigl(\exists c\in\mathcal P(F)\;\;c\text{ is hyperbolic}\bigr)
\quad\Longrightarrow\quad
\#\mathcal P(F)=\infty.
\]
\item Every Finsler metric \(F\) on \(S^2\) has an elliptic prime closed geodesic:
\[
\forall F\;\exists c\in\mathcal P(F)\quad c\text{ is elliptic}.
\]
\end{enumerate}
The second assertion uses the ellipticity convention stated above, subject to the unresolved source qualification. The source also does not define ``distinct'': the first assertion is written using periodic geodesics modulo parameter translation, while the intended alternative of distinct images on \(S^2\) is not explicitly resolved by the supplied excerpt. No additional nondegeneracy or curvature hypothesis is supplied.
\subsection{Short English statement}
The record presents two conjectural claims about Finsler metrics on \(S^2\): one hyperbolic prime closed geodesic forces infinitely many distinct prime closed geodesics, and every such metric has an elliptic prime closed geodesic. The meaning of elliptic at eigenvalues \(\pm1\), and the intended equivalence defining distinct geodesics, require clarification.
\subsection{Sources}
\begin{itemize}
\item Supplied local UnsolvedMath registry record, permanent ID \(800923\), exact source problem code \texttt{catalogue entry}, in the AMR Open Problem Lists collection. The registry labels the record \texttt{open}; this is supplied metadata, not an independently verified mathematical or literature assessment.
\item Burns and Matveev, \emph{Open Problems and Questions about Geodesics}, listed in the supplied record as 2019, Conjecture 3: \url{https://arxiv.org/abs/1308.5417}. The displayed conjecture is attributed to Long; the background also records Long as attribution/context. The second assertion follows the conjecture environment in the supplied excerpt, without a separate attribution label there.
\item Supplied provenance: extraction method \texttt{arxiv\_latex}; source access date \texttt{2026-10-16}. This date records supplied source-access curation and does not indicate independent source verification for this draft.
\end{itemize}
\end{document}
